\chapter{สถาปัตยกรรมการจำลองฟิสิกส์ 3 มิติด้วย Three.js และ PBR Shaders}

การพัฒนาระบบห้องปฏิบัติการเสมือนจริงในปัจจุบัน จำเป็นต้องอาศัยการผสมผสานระหว่างฟิสิกส์เชิงทฤษฎีและคอมพิวเตอร์กราฟิกส์สมรรถนะสูง เพื่อให้ผู้เรียนเกิดประสบการณ์เชิงลึกเสมือนได้สัมผัสอุปกรณ์จริงในห้องปฏิบัติการ

\section{รากฐานระบบเรนเดอร์ WebGL และการบริหารจัดการหน่วยความจำ}

การทำงานของ WebGL บนเว็บเบราว์เซอร์สมัยใหม่ เป็นการสื่อสารโดยตรงกับหน่วยประมวลผลกราฟิกส์ (GPU) ผ่าน Vertex Shaders และ Fragment Shaders โดยมีสถาปัตยกรรมหลักดังแสดงในกรอบแนวคิด

\begin{theoremlab}{หลักการสำคัญ การลดทอนการจัดสรรหน่วยความจำ (Zero Garbage Collection)}
ในแอนิเมชันความเร็ว 60 เฟรมต่อวินาที การสร้างวัตถุใหม่ในลูปเรนเดอร์ เช่น \texttt{new THREE.Vector3()} ส่งผลให้ Garbage Collector ทำงานบ่อยครั้ง เกิดอาการกระตุก (Frame Drop) แนวทางที่ถูกต้องคือการใช้ \textbf{Object Reuse} และคำนวณผ่านตัวแปรระดับโกลบอลที่จัดสรรไว้ล่วงหน้า
\end{theoremlab}

\section{การสร้างพื้นผิว Physically Based Rendering และแสงเงา}

การสร้างแบบจำลองอุปกรณ์วิทยาศาสตร์ เช่น ตัวต้านทาน ขดลวด หรือลูกตุ้มโลหะ ต้องใช้ \texttt{MeshStandardMaterial} หรือ \texttt{MeshPhysicalMaterial} ซึ่งจำลองการสะท้อนของแสงตามหลักทัศนศาสตร์จริง:

\begin{equation}
f_r(\omega_i, \omega_o) = \frac{D(h) F(\omega_i, h) G(\omega_i, \omega_o, h)}{4 (\omega_i \cdot n)(\omega_o \cdot n)}
\end{equation}

โดยที่ $D(h)$ คือการกระจายของพื้นผิวจุลภาค (Microfacet Distribution), $F$ คือพจน์เฟรสเนล (Fresnel Reflectance), และ $G$ คือสัมประสิทธิ์การบดบังของพื้นผิว (Geometric Attenuation)

\section{การจำลองฟิสิกส์เพนดูลัมฮาร์มอนิกและการลดรูปวงจร 6 เหลี่ยม}

ในระบบจำลองการเคลื่อนที่ของลูกตุ้มเพนดูลัมแบบง่าย คาบเวลาการแกว่งที่มุมขนาดเล็ก ($\theta < 10^\circ$) เป็นไปตามสมการ:

\begin{equation}
T = 2\pi \sqrt{\frac{L}{g}}
\end{equation}

ส่วนในวงจรไฟฟ้าแบบโครงข่ายหกเหลี่ยมด้านเท่า (Hexagonal Resistor Network) ที่มีตัวต้านทานขนาด $R$ เท่ากันทั้ง 6 ด้าน เมื่อต่อขั้วขนานระหว่างจุดยอดตรงข้าม ความต้านทานสมมูลรวม $R_{eq}$ สามารถคำนวณได้จากหลักการสมมาตรของศักย์ไฟฟ้า:

\begin{equation}
R_{eq} = \frac{5}{6}R
\end{equation}

เมื่อ $R = 1.0\,\Omega$ จะได้ค่าความต้านทานสมมูลเท่ากับ $0.8333\,\Omega$ ซึ่งตรงกับผลลัพธ์ที่วัดได้จริงจากแอมมิเตอร์และโวลต์มิเตอร์ในแบบจำลอง 3 มิติ
